\documentclass[12pt,a4paper]{article}
\usepackage[margin=2.5cm]{geometry}
\usepackage{amsmath,graphicx,hyperref,natbib,booktabs,xcolor}
\definecolor{ai4sgreen}{HTML}{0D3B2E}
\definecolor{ai4sgold}{HTML}{C8956C}

\title{\textbf{\textcolor{ai4sgold}{AI4S-Journal} Negative Results \& Lessons Learned}}
\author{Author Name\textsuperscript{1*}, Co-Author\textsuperscript{2} \
\textsuperscript{1}Department, Institution, City, Country \
\textsuperscript{2}Department, Institution, City, Country \
\textsuperscript{*}Email: email@institution.edu}
\date{}

\begin{document}
\maketitle

\begin{abstract}
Briefly describe the research question, the approach, the negative or null
result obtained, and the actionable lesson learned. AI4S-Journal reviews
negative results for scientific rigor, not for "success."
\end{abstract}
\noindent\textbf{Keywords:} negative results, lessons learned, reproducibility, [domain]

\section{Introduction}
State the research question and why it is important. Explain why a null or
negative result is scientifically meaningful in this context.

\section{Approach}
\subsection{Hypothesis}
Clearly state the hypothesis that was tested.
\subsection{Experimental Design}
Describe the setup, datasets, models, and evaluation protocol for reproducibility.

\section{Results}
\subsection{Primary Outcome}
Present the negative or null result objectively.
\subsection{Additional Analyses}
Post-hoc explorations of why the hypothesis was not supported.

\section{Lessons Learned}
This is the core contribution of a negative results paper. Clearly articulate:

\begin{enumerate}
    \item \textbf{What was expected, and what happened instead.}
    \item \textbf{Why the result is useful for the community.}
    \item \textbf{What others should do differently.}
    \item \textbf{Recommendations for experimental design going forward.}
\end{enumerate}

\section{Discussion}
Contextualize within the broader literature. If possible, compare with related
negative results.

\section{Conclusion}
\section*{Data and Code Availability}

\bibliographystyle{unsrt}
\begin{thebibliography}{99}
\bibitem{n1} Author, A. (Year). \textit{Title}. Journal.
\end{thebibliography}
\end{document}
